Das
\definitionsverweis {Gradientenfeld}{}{}
zu $h$ ist $\begin{pmatrix}  2x \\ 2y   \end{pmatrix}$, ein Einheitsnormalenfeld ist somit

\mavergleichskettedisp
{\vergleichskette
{ N \begin{pmatrix}  x  \\ y   \end{pmatrix}
}
{ =} { { \frac{   1 }{  \sqrt{x^2+y^2}  } } \begin{pmatrix}  x  \\ y   \end{pmatrix}
}
{ } {
}
{ } {
}
{ } {
}
}
{}{}{.}
Das
\definitionsverweis {totale Differential}{}{}
von $N$ in einem Punkt

\mavergleichskettedisp
{\vergleichskette
{ P
}
{ =} { \begin{pmatrix}  x  \\ y   \end{pmatrix}
}
{ } {
}
{ } {
}
{ } {
}
}
{}{}{}
ist

\mavergleichskettedisp
{\vergleichskette
{  \left(D N \right)_{P}
}
{ =} { \begin{pmatrix}  { \frac{   y^2 }{   \left(  x^2+y^2  \right)^{3/2}  } }   &   - { \frac{   xy }{   \left(  x^2+y^2  \right)^{3/2}  } }   \\    - { \frac{   xy }{   \left(  x^2+y^2  \right)^{3/2}  } }    &  { \frac{   x^2 }{   \left(  x^2+y^2  \right)^{3/2}  } }  \end{pmatrix}
}
{ } {
}
{ } {
}
{ } {
}
}
{}{}{.}
Angewendet auf den tangentialen Vektor
\mathl{\begin{pmatrix}  y  \\ -x   \end{pmatrix}}{} ergibt sich

\mavergleichskettealign
{\vergleichskettealign
{ \left(D N \right)_{P} \begin{pmatrix}  y  \\ -x   \end{pmatrix}
}
{ =} { \begin{pmatrix}  y^2    &   - xy    \\  - xy   &  x^2  \end{pmatrix} \begin{pmatrix}  y  \\ -x   \end{pmatrix}
}
{ =} { \begin{pmatrix}  y \left(  y^2 +x^2  \right)  \\ -x \left(  y^2 +x^2  \right)    \end{pmatrix}
}
{ =} { \begin{pmatrix}  y   \\ -x    \end{pmatrix}
}
{ } {
}
}
{}
{}{.}
Daher ist die Krümmung der Kurve in $P$ nach
[[Ebene differenzierbare Kurve/Faser/Krümmung/Fakt|Fakt]]
gleich $-1$.
\zusatzklammer {Dies passt, da
\mathl{\begin{pmatrix}  y  \\ -x   \end{pmatrix}}{} und das Gradientenfeld die Standardorientierung repräsentieren und
\mathl{\begin{pmatrix}  y  \\ -x   \end{pmatrix}}{} die Tangentialrichtung für die Bewegung mit dem Uhrzeigersinn ist.} {} {}

 [[Kategorie:Latexseite]]
